% Ported from sections/02_results.md (v6, APPROVED 2026-09-18).
% Final-manuscript Sec.~III. Floats: Table 3, Table 4, Figure 2, Figure 3.
\section{Results}
\label{sec:results}

\subsection{The two-phase landscape of the metal film on MgO}
\label{sec:landscape}

Our target is the wetting of the metal film on MgO(001). The property that distinguishes a wetting
from a dewetting film is its flatness, measured by the vertical span of the film,
$\dZ = z(\text{metal})_\text{max} - z(\text{metal})_\text{min}$ over the film's metal atoms (Fe);
$\dZ \approx 0$ is a flat, wetting film and large $\dZ$ a three-dimensional island. We therefore
designed the search to return a map of this two-phase landscape: a GOFEE global optimisation
(Sec.~\ref{sec:bias}) seeded from a flat reference layer, so that the flat basin is
populated from the outset and any lower-lying basin the search finds is populated alongside it. The
result is a sampled distribution over $\dZ$ and energy---an \emph{exploration} density, not a
thermodynamic density of states. Because the sampling is deliberately biased, we use it only for
unweighted structural and energetic comparisons; the number of structures in a basin is a descriptive
attribute of the sampled database, not a physical population.

Across the two models the search returns 1723 structures spanning $\dZ$ from 0.001 to 5.57~\AA,
with both phases populated in each (Table~\ref{tab:phases}). Fig.~\ref{fig:landscape} draws that
landscape: every structure is placed at its own vertical span and at its energy above the model's
lowest, so the figure is the map the search was designed to produce rather than a selection of
representative structures. The points form a continuous band that descends from the flat window into
an island minimum near $\dZ = 3.7$~\AA{} in both models, and the two panels share one energy scale,
so a given height means the same relative energy in either.

\begin{table}[t]
\small
\caption{\label{tab:phases}The two phases in each model ($\dZ \le 1.0$~\AA{} $=$ flat/wetting;
$\dZ > 1.0$~\AA{} $=$ island).}
\begin{tabular*}{\textwidth}{@{\extracolsep{\fill}}lccccc}
\toprule
Model & Structures & Island & Flat & Global-min $\dZ$ (\AA) & Flat-basin min $dE/N$ (\eVA) \\
\midrule
Fe/MgO   & 1180 & 985 & 195 & 3.65 & 0.1888 \\
Fe-B/MgO & 543  & 430 & 113 & 3.77 & 0.1493 \\
\bottomrule
\end{tabular*}
\end{table}

\begin{figure}[t]
\centering
\includegraphics[width=\textwidth]{pes_2_systems.png}
\caption{\label{fig:landscape}The sampled two-phase landscape of the metal film on MgO. Each point is
one structure from the Fe/MgO (left) and Fe-B/MgO (right) searches, plotted as its relative energy
$dE/N$ above the model's own lowest energy against the vertical span of the film, $\dZ$. The shaded
band is the flat window ($\dZ \le 1.0$~\AA); the star marks the model's lowest-energy structure and
the diamond the best member of its flat basin. Both panels use the same energy scale.}
\end{figure}

\subsubsection{The island is the ground state in both models, despite the flat bias} The lowest-energy
structure found is an island in each model, with its global minimum at $\dZ = 3.65$~\AA{} (Fe/MgO)
and 3.77~\AA{} (Fe-B/MgO) (Table~\ref{tab:phases}, column 5)---a three-dimensional cluster rather
than a flat film. This result is \emph{against} the direction of the bias: every search was seeded
from a flat reference layer and, in the early phases, perturbed only at small scale from it, so the
flat configuration had a systematic head start. The search nevertheless left the flat basin and
converged on an island in both models. The flat-wetting state is therefore not an artefact of how the
initial structures were built---it is genuinely higher in energy than the island.

\subsubsection{The flat configuration is a distinct, higher-energy basin} Taking for each model the
lowest-energy structure with $\dZ \le 1.0$~\AA{} isolates the flat basin's best member. In both
models it lies above the global minimum---by 0.1888~\eVA{} for Fe/MgO and 0.1493~\eVA{} for
Fe-B/MgO (Table~\ref{tab:phases}, column 6). The flat film is thus a separate basin of the
landscape, not merely the high-$\dZ$ tail of a single minimum.

\subsubsection{The branches are families, not single structures} The $\dZ$ values are continuous rather
than grouped into discrete island heights: in each model the island branch spans $\dZ \approx
1$--$5.6$~\AA{} with a median of 3.02~\AA{} (Fe/MgO) and 2.58~\AA{} (Fe-B/MgO), and the flat branch
spans the whole window below 1~\AA. The low-energy members of each branch are likewise not a single
repeated structure: comparing the lowest-energy structures of each branch by their radial/angular
fingerprint shows that each branch's low-energy set splits into two distinct structural
motifs, and that a motif recurs across independent searches---for Fe/MgO, the five lowest flat
structures come from five different searches and fall into two motifs, four of them in the dominant
one.
The fingerprint is computed on the whole template-plus-film structure, so the fixed MgO template
compresses the pair distances and the motif count is a lower bound on the structural
diversity present; a film-resolved descriptor would be sharper. Pairwise distances within the
low-energy set of a branch are 0.41--0.44 of that branch's random-pair distance scale, i.e.\ the
low-energy structures are consistently more compact than the branch as a whole. The count is also
limited by how many searches contribute to the low-energy set: the Fe-B/MgO flat set within
0.02~\eVA{} of its minimum comes from two searches and its island set from three, so those counts
rest on fewer independent searches than the corresponding Fe/MgO numbers (five and four).

\subsubsection{Experimental correspondence} The films in this study are about one monolayer thick
on average---25 metal atoms per 25 substrate sites in the $5\times5$ cell, the boron-containing film
carrying three additional boron atoms---and that is the regime in which the growth mode of Fe on
MgO(001) has been characterised experimentally. At room temperature the observed growth is
three-dimensional: He-atom scattering shows 3D metal island growth of Fe on MgO(001), with
quantified island densities, sizes and shapes, suppressed only when the film is deposited at low
temperature (140~K), where a monolayer almost completely covers the substrate \cite{fahsold2000};
scanning tunnelling microscopy finds that sub-nanometre Fe grows three-dimensionally on MgO,
with island coalescence between 3.5 and 6.5 monolayers and a two-dimensional growth mode reported
only above $\approx$6.5 monolayers \cite{torelli2009}; and grazing-incidence small-angle X-ray
scattering on five monolayers evaporated at room temperature likewise identifies Volmer--Weber
growth with spherical islands \cite{reitinger2007}. A flat, continuous film is obtained instead by
low-temperature deposition \cite{fahsold2000} or by slow deposition onto a cleaved, oxygen-annealed
crystal, where the first monolayer grows pseudomorphically and layer by layer \cite{urano1988}. Our
two basins therefore map onto the two outcomes that are experimentally accessible at this
coverage---the flat basin onto the pseudomorphic monolayer and the island basin onto the 3D
clusters---and the energy ordering we find places the flat, wetting configuration above the island.
This is a consistency check rather than a prediction: the observed growth mode is also set by
kinetics, so the correspondence is between structural configurations, not between energies. Because
the growth mode turns two-dimensional only above $\approx$6.5 monolayers \cite{torelli2009}, the
comparison is specific to the monolayer regime and should not be extrapolated to thick films.

\subsection{The flat (wetting) phase with and without boron}
\label{sec:flat}

With the two phases established, we compare how adding boron to the Fe film shifts the flat basin.
Table~\ref{tab:flatbasin} gives the flat-basin minimum $dE/N$ for the two models, i.e.\ the same host
with and without the metalloid, and Fig.~\ref{fig:flat} places the two values side by side with the
change between them marked.

\begin{table}[t]
\small
\caption{\label{tab:flatbasin}Flat-basin minimum $dE/N$ (\eVA) for the same Fe host with and
without boron.}
\begin{tabular*}{\textwidth}{@{\extracolsep{\fill}}lcc}
\toprule
Film & Flat-basin min $dE/N$ (\eVA) & B effect \\
\midrule
Fe   & 0.1888 & --- \\
Fe-B & 0.1493 & \textbf{$-0.040$} \\
\bottomrule
\end{tabular*}
\end{table}

\begin{figure}[t]
\centering
\includegraphics[width=\textwidth]{flat_state_summary.png}
\caption{\label{fig:flat}Flat-basin minimum relative energy for the same Fe host with and without
boron. A bar is the lowest energy found among the structures of the flat window ($\dZ \le 1.0$~\AA),
in \eVA{} above the model's own lowest energy, and the bracket gives the shift produced by boron.}
\end{figure}

\subsubsection{Boron lowers the flat-state energy} Adding boron lowers the flat basin's best member by
0.040~\eVA---from 0.1888 to 0.1493~\eVA---a 21\,\% reduction of the flat--island separation.
The effect is measured across 13 independent Fe/MgO searches and 6 independent Fe-B/MgO searches; its
search-level significance is given in the uncertainty paragraph below.

\subsubsection{Boron does not bond to the MgO interface} Restricting to the low-energy window
($dE/N \le 0.05$~\eVA), only 1 of the 72 Fe-B/MgO structures in that window has any boron atom within
2.6~\AA{} of an oxygen of the substrate; the maximum boron--oxygen contact fraction in the window is
0.33. Outside the window the picture is different---101 of the 471 remaining Fe-B/MgO structures do
have a boron--oxygen contact---which is why the statement is restricted to the low-energy window
rather than made over all sampled structures. Boron therefore remains inside the metal film rather
than wetting the oxide, so its effect on the flat phase is film-internal rather than a change in
interfacial bonding.

\subsubsection{Uncertainty from the choice of search} Each model was searched from several independent
random seeds (13 completed searches for Fe/MgO, 6 for Fe-B/MgO), and structures from one search are
correlated with each other, so the independent unit is the search, not the structure. Taking each
search's own flat-basin minimum, the per-search minimum falls from 0.2171~\eVA{} without boron to
0.1882~\eVA{} with it on average (medians 0.2185 and 0.1834), a median shift of
0.0351~\eVA{} with 94\,\% of resampled replicates agreeing in sign---the effect is not carried by a
single favourable search, and the best-structure difference of Table~\ref{tab:flatbasin}
($-0.040$~\eVA) is the extreme of the same distribution rather than a lone outlier. A two-sided
permutation test over the 19 searches gives $p = 0.0444$, significant at the 5\,\% level. In
words, $p$ is the probability that a random reshuffling of the 19 per-search minima into two groups
of 13 and 6 would, by chance alone, produce a shift in the flat-state energy at least as large as the
one observed---so $p = 0.0444$ means such a shift arises by chance in only about 4.4\,\% of
reshufflings, i.e.\ the observed lowering of the flat state is unlikely to be a chance outcome of
which searches happened to contain boron. The test is an exact enumeration of all 27\,132
partitions of the pooled search minima, so this value is exact for the data and its resolution floor
is $p = 3.7\times10^{-5}$---more than three orders of magnitude below the value reported, i.e.\ the
conclusion is not an artefact of too few searches.

\subsection{The island (dewetting) phase and the separation between phases}
\label{sec:island}

Adding boron leaves the \emph{island} phase qualitatively similar. In both models the lowest-energy
structure is an island---at $\dZ = 3.65$~\AA{} (Fe/MgO) and 3.77~\AA{} (Fe-B/MgO)
(Table~\ref{tab:phases})---and the island branch spans the same broad range of heights in both
($\dZ \approx 1$--$5.6$~\AA, median 3.02~\AA{} without boron and 2.58~\AA{} with it). Boron does not
change what the dewetted state looks like; it changes how far above it the flat, wetting state sits.

Because each model is referenced to its own lowest energy, the quantity that compares the two phases
is the separation between them---the flat-basin minimum of Table~\ref{tab:flatbasin}---and
that separation is what boron reduces. The present data do not resolve \emph{how} boron reduces it: a
lower flat-basin energy could equally arise from stabilising the flat phase or from destabilising the
island, and distinguishing the two would require the island phase to be referenced to an absolute
(rather than per-model) energy scale, or a converged treatment of both basins. We therefore state the
effect as a change in the flat--island separation, not as the stabilisation of either phase
individually.

\subsection{Summary of the results}
\label{sec:summary}

The biased GOFEE search, seeded from a flat reference, maps both the flat (wetting) and island
(dewetting) phases in both models, over $\dZ = 0.001$--$5.57$~\AA. The island is the ground state in
both models ($\dZ = 3.65$ and 3.77~\AA), even though the search was biased toward the flat
configuration---the flat bias did not manufacture the result---and the flat configuration is a
distinct, higher-energy basin, 0.1888~\eVA{} above the island for Fe/MgO and 0.1493~\eVA{} for
Fe-B/MgO, in both cases a separate basin rather than the high-$\dZ$ tail of a single minimum. Boron
lowers the flat--island separation by 0.040~\eVA{} (13 vs 6 completed searches, exact two-sided
permutation $p = 0.0444$) without displacing the island as the ground state, and it does not bond to
the MgO interface---it stays within the film, so its effect on the flat phase is film-internal.

The low-energy structures of each branch form a small set of recurring motifs rather than one
repeated structure, and $\dZ$ is continuous: the two phases are families, and the energetics above
are compared unweighted, since the sampled densities are biased and not physical.

The mechanistic origin of these trends (electronic structure, interface registry) and the robustness
of the method to its parameters are presented in the Supplementary Material.
